% preamble-exam-zh.tex — 共享 preamble
% 用于所有 exam-zh 模板
%
% 样式策略：
%   屏幕 → 语义彩色（蓝=关键想法，红=易错，绿=路线，灰=提示/教师备注）
%   打印 → 自动降级为黑白（通过 print mode 或 monochrome 选项）

\documentclass{exam-zh}

% ---------- 基础配置 ----------
\examsetup{
  page/size = a4paper,
  page/show-head = false,
  page/show-foot = false,
  sealline/show = false,
  font = times,
  math-font = xits,
}

\usepackage{tcolorbox}
\usepackage{needspace}
\tcbuselibrary{skins,breakable}

% ---------- 打印/屏幕颜色切换 ----------
% 用法：\documentclass[monochrome]{exam-zh} 即可切换为纯黑白
% 注意：不要使用 \if@xxx 放在 makeatother 外部；这里改成普通条件名。
\newif\ifeduprintcolor
\eduprintcolortrue

\makeatletter
\@ifclasswith{exam-zh}{monochrome}{\eduprintcolorfalse}{}
\makeatother

% ---------- 语义颜色定义 ----------
% 屏幕：有色彩区分；打印：降级为灰度
\ifeduprintcolor
  % --- 屏幕彩色模式 ---
  \definecolor{edu-blue}{RGB}{47, 82, 122}       % 关键想法
  \definecolor{edu-blue-bg}{RGB}{243, 246, 252}
  \definecolor{edu-red}{RGB}{180, 40, 40}         % 易错提醒
  \definecolor{edu-red-bg}{RGB}{255, 245, 240}
  \definecolor{edu-green}{RGB}{34, 100, 60}       % 路线图
  \definecolor{edu-green-bg}{RGB}{242, 250, 245}
  \definecolor{edu-orange}{RGB}{160, 90, 0}       % 训练目标
  \definecolor{edu-orange-bg}{RGB}{255, 248, 235}
  \definecolor{edu-gray}{RGB}{100, 100, 100}      % 提示/教师备注
  \definecolor{edu-gray-bg}{RGB}{248, 248, 248}
  \definecolor{edu-purple}{RGB}{90, 50, 120}      % 思考题
  \definecolor{edu-purple-bg}{RGB}{248, 244, 252}
\else
  % --- 打印黑白模式 ---
  \definecolor{edu-blue}{gray}{0.15}
  \definecolor{edu-blue-bg}{gray}{0.96}
  \definecolor{edu-red}{gray}{0.25}
  \definecolor{edu-red-bg}{gray}{0.97}
  \definecolor{edu-green}{gray}{0.20}
  \definecolor{edu-green-bg}{gray}{0.97}
  \definecolor{edu-orange}{gray}{0.25}
  \definecolor{edu-orange-bg}{gray}{0.98}
  \definecolor{edu-gray}{gray}{0.40}
  \definecolor{edu-gray-bg}{gray}{0.97}
  \definecolor{edu-purple}{gray}{0.30}
  \definecolor{edu-purple-bg}{gray}{0.98}
\fi

% ---------- 自定义教学环境 ----------
% 共性：enhanced + breakable（跨页不断裂）

% 原题卡片 — 强边框，突出显示
\newtcolorbox{problemcard}{
  enhanced, breakable,
  colback=edu-blue-bg,
  colframe=edu-blue,
  boxrule=1.2pt,
  arc=0pt,
  left=3mm, right=3mm, top=3mm, bottom=3mm,
}

% 解题路线图 — 左边线 + 浅底
\newtcolorbox{routemap}{
  enhanced, breakable,
  colback=edu-green-bg,
  colframe=edu-green!30,
  boxrule=0pt,
  borderline west={2.5pt}{0pt}{edu-green},
  left=3mm, right=3mm, top=3mm, bottom=3mm,
}

% 关键想法 — 蓝色左边线
\newtcolorbox{keyideabox}{
  enhanced, breakable,
  colback=edu-blue-bg,
  colframe=edu-blue!30,
  boxrule=0pt,
  borderline west={2.5pt}{0pt}{edu-blue},
  left=3mm, right=3mm, top=3mm, bottom=3mm,
}

% 易错提醒 — 红色左边线
\newtcolorbox{mistakebox}{
  enhanced, breakable,
  colback=edu-red-bg,
  colframe=edu-red!20,
  boxrule=0pt,
  borderline west={3pt}{0pt}{edu-red},
  left=3mm, right=3mm, top=2.5mm, bottom=2.5mm,
}

% 提示 — 灰色虚线左边线
\newtcolorbox{hintbox}{
  enhanced, breakable,
  colback=edu-gray-bg,
  colframe=edu-gray!20,
  boxrule=0pt,
  borderline west={2pt}{0pt}{edu-gray, dashed},
  left=3mm, right=3mm, top=2mm, bottom=2mm,
}

% 思考题 — 紫色虚线左边线
\newtcolorbox{thinkbox}{
  enhanced, breakable,
  colback=edu-purple-bg,
  colframe=edu-purple!20,
  boxrule=0pt,
  borderline west={2pt}{0pt}{edu-purple, dashed},
  left=2.5mm, right=2.5mm, top=2.5mm, bottom=2.5mm,
}

% 教师备注 — 灰色左边线，小字
\newtcolorbox{teachernote}{
  enhanced, breakable,
  colback=edu-gray-bg,
  colframe=edu-gray!15,
  boxrule=0pt,
  borderline west={2pt}{0pt}{edu-gray!60},
  left=2.5mm, right=2.5mm, top=2.5mm, bottom=2.5mm,
  fontupper=\small,
  colupper=black!60,
}

% 训练目标 — 橙色左边线，小字
\newtcolorbox{traininggoal}{
  enhanced, breakable,
  colback=edu-orange-bg,
  colframe=edu-orange!15,
  boxrule=0pt,
  borderline west={1.5pt}{0pt}{edu-orange},
  left=2mm, right=2mm, top=1mm, bottom=1mm,
  fontupper=\small,
}

% 答题区 — 无边框，纯留白
\newcommand{\answerarea}[1][40mm]{%
  \vspace{2mm}%
  \begin{tcolorbox}[
    enhanced, breakable,
    colback=white,
    colframe=white,
    boxrule=0pt,
    height=#1,
    valign=top,
    bottom=0mm,
    after skip=0mm,
  ]
  \end{tcolorbox}%
}

% ---------- 字体设置 ----------
% exam-zh (ctexbook) already sets CJK fonts via ctex.
% Only override if specific fonts are needed.
% macOS: Songti SC, STHeiti, STFangsong
% Linux: SimSun, SimHei, FangSong

% ---------- 页面设置 ----------
% exam-zh (ctexbook) already loads geometry.
% Override margins if needed:
% \geometry{top=16mm, bottom=16mm, left=14mm, right=14mm}

\examsetup{
  question/show-answer = false,
  fillin/show-answer = false,
  solution/show-solution = hide,
}

% ---------- 教师版解答样式 ----------
\newtcolorbox{solutionblock}{
  enhanced, breakable,
  colback=edu-blue-bg,
  colframe=edu-blue!25,
  boxrule=0pt,
  borderline west={2pt}{0pt}{edu-blue!50},
  arc=0pt,
  left=3mm, right=3mm, top=2mm, bottom=2mm,
  before skip=2mm,
  after skip=3mm,
  fontupper=\small,
}

% 解答步骤编号
\newcounter{solstep}
\newcommand{\solstep}[2]{%
  \stepcounter{solstep}%
  \par\medskip
  \noindent\hangindent=6mm
  {\color{edu-blue}\bfseries\textcircled{\scriptsize\thesolstep}\enspace #1}\enspace #2\par
}

\begin{document}

\section*{等腰三角形与直角三角形存在性 · 专题练习}

\section*{一、选择题（每题 3 分，共 18 分）}


\begin{question}[points=3]
  在 $x$ 轴上存在点 $P$ 使 $\triangle PAB$ 为等腰三角形。已知 $A(0,3)$，$B(4,0)$，则下列结论正确的是
  \begin{choices}
    \item 只有 $2$ 个符合条件的 $P$ 点
    \item 一定有 $3$ 个符合条件的 $P$ 点
    \item 最多有 $6$ 个符合条件的 $P$ 点（三等腰，每种最多 $2$ 个解）
    \item 条件不足，无法判断
  \end{choices}
\end{question}



\begin{question}[points=3]
  已知 $A(-1,2)$，$B(3,2)$，$P$ 在 $y$ 轴上，$\triangle ABP$ 为直角三角形。当 $\angle P = 90°$ 时，$P$ 的纵坐标满足的方程是
  \begin{choices}
    \item $y^2 - 4y + 1 = 0$
    \item $y^2 - 4y + 3 = 0$
    \item $y^2 + 4y - 3 = 0$
    \item $y^2 + 4y + 3 = 0$
  \end{choices}
\end{question}



\begin{question}[points=3]
  $\triangle ABC$ 中 $A(0,0)$，$B(6,0)$，$C(0,8)$。$P$ 在 $x$ 轴上使 $\triangle PAC$ 为等腰三角形。下列哪种情况在 $x$ 轴上一定有解？
  \begin{choices}
    \item $PA = PC$（$P$ 在 $AC$ 的垂直平分线上）
    \item $AP = AC$
    \item $CP = CA$
    \item 三种情况在 $x$ 轴上都有解
  \end{choices}
\end{question}



\begin{question}[points=3]
  直角 $\triangle ABC$ 中 $\angle C = 90°$，$AC=3$，$BC=4$。$P$ 在 $AB$ 上运动，使 $\triangle PBC$ 为等腰三角形。下列哪种情况一定有解？
  \begin{choices}
    \item $PB=PC$（$P$ 在 $BC$ 的垂直平分线上）
    \item $BP=BC=4$
    \item $CP=CB=4$
    \item 以上都不一定
  \end{choices}
\end{question}



\begin{question}[points=3]
  已知 $A(2,0)$，$B(0,4)$。$P$ 在 $x$ 轴上使 $\triangle ABP$ 为直角三角形。$\angle A=90°$ 时，$P$ 的坐标为
  \begin{choices}
    \item $P(2,0)$（即 $P=A$，退化）
    \item $P(2,0)$
    \item 不存在符合条件的 $P$
    \item $P$ 有两个解
  \end{choices}
\end{question}



\begin{question}[points=3]
  下列关于等腰三角形存在性问题的说法，正确的是
  \begin{choices}
    \item 等腰三角形一定有 $3$ 个解（三等腰各一个）
    \item 求出解后不需要检验三点是否共线
    \item 当动点与某个已知顶点重合时，三角形退化为线段，该解需要排除
    \item 直角三角形只需要讨论两种情况（直角在动点和直角在一个定点）
  \end{choices}
\end{question}


\section*{二、填空题（每题 3 分，共 18 分）}


\begin{question}[points=3]
  $A(0,0)$，$B(4,0)$，$C(2,3)$。$P$ 在 $x$ 轴上使 $\triangle PBC$ 为等腰三角形。当 $BP=BC$ 时，$P$ 的横坐标为 $\underline{\quad\quad}$。
  \fillin
\end{question}



\begin{question}[points=3]
  $A(1,0)$，$B(0,2)$。$P$ 在 $y$ 轴上使 $\triangle ABP$ 为直角三角形。$\angle A=90°$ 时，$P$ 的纵坐标为 $\underline{\quad\quad}$。
  \fillin
\end{question}



\begin{question}[points=3]
  正方形 $ABCD$ 中 $A(0,0)$，$B(4,0)$，$C(4,4)$，$D(0,4)$。$P$ 在直线 $y=x$ 上使 $\triangle ABP$ 为等腰三角形。当 $AP=AB$ 时，$P$ 到原点的距离为 $\underline{\quad\quad}$。
  \fillin
\end{question}



\begin{question}[points=3]
  已知 $A(-2,0)$，$B(2,0)$，$P$ 在直线 $y=x+1$ 上使 $\triangle ABP$ 为直角三角形。$\angle P=90°$ 时，$P$ 的坐标满足的方程为 $\underline{\quad\quad}$（化为一般式）。
  \fillin
\end{question}



\begin{question}[points=3]
  矩形 $ABCD$ 中 $A(0,0)$，$B(6,0)$，$C(6,4)$，$D(0,4)$。$P$ 在 $BC$ 上使 $\triangle APD$ 为等腰三角形。当 $AP=DP$ 时，$P$ 的坐标为 $\underline{\quad\quad}$。
  \fillin
\end{question}



\begin{question}[points=3]
  $\triangle ABC$ 中 $A(0,0)$，$B(6,0)$，$C(3,4)$。$P$ 在 $y$ 轴上使 $\triangle PBC$ 为等腰三角形。$PC=BC$ 时，$P$ 的坐标为 $\underline{\quad\quad}$。
  \fillin
\end{question}


\section*{三、解答题（每题 10 分，共 30 分）}

\needspace{8\baselineskip}

\begin{problem}[points=10]
  直线 $y=-\dfrac{4}{3}x+4$ 与 $x$ 轴、$y$ 轴分别交于 $A$、$B$ 两点。$P$ 是 $x$ 轴上的动点，$\triangle ABP$ 为等腰三角形。

（1）求 $A$、$B$ 的坐标和 $AB$ 的长；

（2）求所有符合条件的 $P$ 的坐标。
  \answerarea[80mm]
\end{problem}


\needspace{8\baselineskip}

\begin{problem}[points=10]
  正方形 $OABC$ 边长 $4$，$O$ 在原点，$A$ 在 $x$ 轴正方向，$C$ 在 $y$ 轴正方向。$D(2,0)$。$P$ 从 $O$ 沿 $O \to A \to B$ 运动（$1$ 单位/秒）。

是否存在 $P$ 使 $\triangle CDP$ 为等腰三角形？若存在，求所有 $P$ 的坐标。
  \answerarea[80mm]
\end{problem}


\needspace{8\baselineskip}

\begin{problem}[points=10]
  $\triangle ABC$ 中 $AB=AC=5$，$BC=6$。$D$ 是 $AB$ 上一点，$DE \parallel BC$ 交 $AC$ 于 $E$。以 $DE$ 为边在 $DE$ 下方作正方形 $DEFG$（$G$ 在 $\triangle ABC$ 内部）。

当 $\triangle BDG$ 为等腰三角形时，求 $AD$ 的长。
  \answerarea[80mm]
\end{problem}


\clearpage
\section*{参考答案}


\begin{keyideabox}
  \textbf{选择题}
  1. C

2. A

3. B

4. B

5. C

6. C
\end{keyideabox}



\begin{keyideabox}
  \textbf{填空题}
  1. $x=4\pm\sqrt{13}$

2. $y=-\dfrac{1}{2}$

3. $4$

4. $2x^2+2x-3=0$

5. $P(6,2)$

6. $P(0,0)$ 或 $P(0,8)$
\end{keyideabox}



\begin{keyideabox}
  \textbf{解答题}
  第 1 题：$P\left(\dfrac{3}{2},0\right)$，$P(8,0)$，$P(-2,0)$，$P(-3,0)$，共 $4$ 个解。

第 2 题：需完整分段计算（与例题12结构类似，共约 $4$ 个解）。

第 3 题：$AD = \dfrac{20}{7}$ 或 $AD = \dfrac{25}{11}$ 或 $AD = \dfrac{125}{73}$。
\end{keyideabox}



\end{document}
